% !PRINTED 396
\weberpage{396}
Thus each of the complexes 6), 7) is azygetic with the complex 1),
and hence one and only one of the two double tangents
$\varrho, \varrho'$ must occur in 6); if this is $\varrho$, then $\varrho'$ is contained in
7). If $\varrho\sigma$ is a pair of 6), then $\sigma$ is different from its
conjugate $\sigma'$, because otherwise $\sigma\varrho'$ would have to occur in 7) and consequently
$\varrho\varrho'$ in 5), which is not the case. Thus, if $\tau, \tau'$ are two further conjugate imaginary double tangents,
we have the complexes:
\[
\begin{array}{rl}
6) & x\xi,\ p\eta,\ r\zeta,\ \varrho\sigma,\ \sigma'\tau \\
7) & x\xi',\ p\eta',\ r\zeta',\ \varrho'\sigma',\ \sigma\tau'.
\end{array}
\]
If $t\lambda$ is the last pair of the complex 6), then the pair
$t'\lambda'$ occurs in 7), and these two pairs must contain a common
element. But since $\lambda$ cannot be equal to $t'$, because otherwise
both pairs would be identical, we must have $t = t'$ (or $\lambda = \lambda'$, which
is not essentially different); thus $t$ is real, and we
obtain:
\[
\begin{array}{rl}
6) & x\xi,\ p\eta,\ r\zeta,\ \varrho\sigma,\ \sigma'\tau,\ t\lambda \\
7) & x\xi',\ p\eta',\ r\zeta',\ \varrho'\sigma',\ \sigma\tau',\ t\lambda'.
\end{array}
\]
We now return to the syzygetic complex triple 1), 2), 3).
Since from 6) we conclude that $x\,p\,t$ and $x\,r\,t$ are azygetic,
$t$ can occur neither in 2) nor in 3), and must therefore be
contained in 1). But since 1) is a real complex, $t$ must appear in it
paired with a real double tangent $u$, and the
complex 1) becomes
\[
x\,y,\ p\,q,\ r\,s,\ t\,u,
\]
whereby our Theorem 3. is proved.

\gesperrt{4. If five real pairs occur in a complex, then
all the double tangents are real.}

Let us start from a complex with five real pairs
\[
1)\quad x_1y_1,\ x_2y_2,\ x_3y_3,\ x_4y_4,\ x_5y_5,
\]
and assume first that besides the last pair of this
complex there exists an imaginary double tangent $\xi$; then
a conjugate double tangent $\xi'$ is also present, and the real complex 2) determined by the
pair $\xi\xi'$ is syzygetic with 1) [because
$\xi$ and $\xi'$ do not occur in 1)]. But the four common
elements of 1) and 2) must be real
%CONT>>
